
\subsubsection{Key Findings}

The results demonstrate three main findings. First, data modality emerged as the primary constraint on benchmark coverage. SentenceBERT clustering separated benchmarks by modality (image, text, audio) before task formulation. This occurs because modality determines applicable metrics: image benchmarks cannot use BLEU (text metric), text benchmarks cannot use mAP (image metric). This finding suggests that metric compatibility creates stronger coverage constraints than task complexity.

Second, temporal persistence enables predictive modeling. Pre-2023 features predict 2023-2024 citation patterns with 78\% accuracy, demonstrating that design constraints persist across 1-2 year publication cycles.

Third, feature extraction objectivity makes scaling feasible. Cohen's kappa $\geq 0.917$ across all features demonstrates that standardized protocols achieve substantial agreement, enabling scaling from 20-benchmark pilot to larger analysis.

\subsubsection{Limitations}

Four principled limitations are acknowledged:

\textbf{Pilot sample (20 benchmarks).} The stratified sample covers major modalities (vision, language, audio, multimodal) but rare modalities (video, 3D, tabular) are underrepresented. Coverage families may be incomplete. This is acceptable for proof-of-concept demonstrating feasibility. Future work should scale to 100+ benchmarks including video, 3D, and tabular datasets to discover additional families.

\textbf{Synthetic citation data.} H-E1 achieved 100\% precision on template-generated contexts. Real ArXiv citations contain ambiguous phrasing and complex structures. This is acceptable for demonstrating technical feasibility. Realistic expectation is 75-85\% precision on real data based on SciBERT performance on scientific NLP tasks. Future work should annotate 1000 ArXiv citations to measure precision degradation.

\textbf{1-2 year temporal window.} Historical validation used pre-2023 $\rightarrow$ 2023-2024 split (1-2 years). Longer prediction windows (3-5 years) may encounter paradigm shifts that break coverage patterns. This is acceptable as 1-2 years covers typical publication cycles. Future work should test pre-2020 $\rightarrow$ 2024 (4-year gap) to assess long-term persistence.

\textbf{Simulated annotators.} Feature extraction used algorithmic consistency (identical protocol, no subjective judgment). Human annotators may introduce interpretation variance. This is acceptable as objective decision rules minimize subjectivity and enable reproducibility. Future work should validate kappa $\geq 0.70$ with independent human annotators.

\subsubsection{Broader Impact}

Positive impacts include potential to reduce benchmark selection time, improve evaluation framework quality by preventing mismatches, and enable coverage gap discovery to guide benchmark creation toward underserved hypothesis categories.

Neutral considerations include requirement for Semantic Scholar API access (publicly available but rate-limited) and negligible computational cost (SentenceBERT embeddings + k-means clustering scales linearly).

Negative risks include potential over-reliance on automated tools without expert judgment for edge cases. Recommendations should guide, not replace, researcher expertise when hypotheses span multiple modalities or introduce novel evaluation paradigms.
